\chapter*{Introduction} \label{I.0}

Nous pr\'esentons ici, sous forme r\'evis\'ee et compl\'et\'ee, une r\'e\'edition par photo-offset du deuxi\`eme S\'EMINAIRE DE G\'{E}OM\'{E}TRIE ALG\'EBRIQUE de l'INSTITUT des HAUTES \'ETUDES SCIENTIFIQUES tenu en 1962 (mim\'eographi\'e).

Le lecteur se reportera \`a l'Introduction au premier de ces S\'eminaires (cit\'e \SGA 1 par la suite) pour les buts que poursuivent ces s\'eminaires, et leurs relations avec les \'EL\'EMENTS DE G\'{E}OM\'{E}TRIE ALG\'EBRIQUE.

Le texte des expos\'es~\ref{I} \`a~\ref{XI} a \'et\'e r\'edig\'e \`a mesure, d'apr\`es mes expos\'es oraux et notes manuscrites, par un groupe d'auditeurs, comprenant I.\, GIROGUITTI, J.\, GIRAUD, Mlle M.\, JAFFE (devenue Mme M.\, HAKIM) et A.\, LAUDAL. Ces notes \`a l'origine \'etaient consid\'er\'ees comme devant \^etre provisoires et \`a circulation tr\`es limit\'ee, en attendant leur absorption par les EGA (absorption devenue maintenant pour le moins probl\'ematique, tout comme pour les autres parties des SGA). Comme il \'etait dit dans l'avertissement \`a l'\'edition primitive, ce caract\`ere \og confidentiel\fg des notes devait excuser certaines \og faiblesses de style\fg, sans doute plus manifestes dans le pr\'esent S\'eminaire \SGA 2 que dans les autres. J'ai essay\'e dans la mesure du possible d'y obvier dans la pr\'esente r\'e\'edition, par une r\'evision relativement serr\'ee du texte initial. J'ai notamment harmonis\'e les syst\`emes de num\'erotation des \'enonc\'es, employ\'es dans les diff\'erents expos\'es, en introduisant partout le m\^eme syst\`eme d\'ecimal, d\'ej\`a utilis\'e dans la plupart des expos\'es primitifs du pr\'esent \SGA 2, ainsi que dans toutes les autres parties des SGA. Cela m'a amen\'e en particulier \`a revoir enti\`erement la num\'erotation des \'enonc\'es des expos\'es~\ref{III} \`a~\ref{VIII}, (et par cons\'equent, des r\'ef\'erences auxdits expos\'es)\footnote{Il va sans dire que toutes les r\'ef\'erences \`a \SGA 2 qui figurent dans les parties des SGA publi\'ees dans la SERIES IN PURE MATHEMATICS se rapporteront au pr\'esent volume, et non \`a l'\'edition primitive de \SGA 2 !}. J'ai essay\'e \'egalement d'extirper du texte primitif les principales erreurs de dactylographie ou de syntaxe (qui \'etaient nombreuses et g\^enantes). De plus, Mme M.\, HAKIM a bien voulu se charger de r\'e\'ecrire l'expos\'e~\ref{IV} dans un style moins t\'el\'egraphique que l'expos\'e initial. Comme dans les autres r\'e\'editions des SGA, j'ai \'egalement ajout\'e un certain nombre de notes de bas de page, soit pour donner des r\'ef\'erences suppl\'ementaires, soit pour signaler l'\'etat d'une question pour laquelle des progr\`es ont \'et\'e faits depuis la r\'edaction du texte primitif. Enfin, ce S\'eminaire a \'et\'e augment\'e d'un nouvel expos\'e, a savoir l'expos\'e~\ref{XIV}, r\'edig\'e par Mme. MICH\`ELE RAYNAUD en 1967, qui reprend et compl\`ete des suggestions contenues dans les \og Commentaires \`a l'Expos\'e~\ref{XIII}\fg (\ref{XIII}~\ref{XIII.6}) (r\'edig\'es en Mars~1963). Cet expos\'e reprend les th\'eor\`emes du type Lefschetz du point de vue de la cohomologie \'etale, en utilisant les r\'esultats sur la cohomologie \'etale expos\'es dans \SGA 4 et \SGA 5 (\`a para\^itre dans cette m\^eme collection SERIES IN PURE MATHEMATICS); il est donc \`a ce titre de nature moins \og \'el\'ementaire\fg que les autres expos\'es du pr\'esent volume, qui n'utilisent gu\`ere plus que la substance des chapitres~I \`a~III des EGA. Voici une esquisse du contenu du pr\'esent volume. L'expos\'e~\ref{I} contient le sorite de la \og cohomologie \`a support dans $Y$\fg $H^*_Y(X, F)$, o\`u $Y$ est un ferm\'e d'un espace~$X$, cohomologie qui peut s'interpr\'eter comme une cohomologie de~$X$ \emph{modulo} l'ouvert $X-Y$, et qui est l'aboutissement d'une fort utile \og \emph{suite spectrale de passage du local au global}\fg \ref{I}~\ref{I.2.6}, faisant intervenir des \emph{faisceaux} de cohomologie \og \`a support dans~$Y$\fg $\underline{H}^*_Y(F)$. Ce formalisme peut dans de nombreuses questions jouer un r\^{o}le de \og localisation\fg analogue \`a celui jou\'e par la consid\'eration de voisinages \og tubulaires\fg de $Y$ en g\'eom\'etrie diff\'erentielle. L'expos\'e~\ref{II} \'etudie les notions pr\'ec\'edentes dans le cas des faisceaux quasi-coh\'erents sur les pr\'esch\'emas, l'expos\'e~\ref{III} donne leur relation avec la notion classique de \emph{profondeur} (\ref{III}~\ref{III.3.3}).

Les expos\'es~\ref{IV} et~\ref{V} donnent des notions de \emph{dualité locale}, qu'on peut comparer au th\'eor\`eme de dualit\'e projective de Serre (\ref{XII}~\ref{XII.1.1}); signalons que ces deux types de th\'eor\`emes de dualit\'e sont g\'en\'eralis\'es de fa\c con substantielle dans le s\'eminaire de HARTSHORNE (cit\'e en note de bas de page \`a la fin de \Exp \ref{IV})\refstepcounter{toto}

Les expos\'es~\ref{VI} et~\ref{VII} donnent des notions techniques faciles, utilis\'ees dans l'expos\'e~\ref{VIII} pour prouver le \emph{théorème de finitude} (\ref{VIII}~\ref{VIII.2.3}), donnant des conditions n\'ecessaires et suffisantes, pour un faisceau coh\'erent $F$ sur un sch\'ema noeth\'erien $X$, pour que les faisceaux de cohomologie locale $\underline{H}^i_Y(F)$ soient coh\'erents pour $i\leq n$ (ou ce qui revient au m\^eme, pour que les faisceaux $\R^if_*(F|X-Y)$ soient coh\'erents pour $i\leq n-1$, o\`u $f\colon X-Y\to X$ est l'inclusion). Ce th\'eor\`eme est un des r\'esultats techniques centraux du S\'eminaire, et nous montrons dans l'expos\'e~\ref{IX} comment un th\'eor\`eme de cette nature peut \^etre utilis\'e, pour \'etablir un \og th\'eor\`eme de comparaison\fg et un \og th\'eor\`eme d'existence\fg en g\'eom\'etrie formelle, en calquant et g\'en\'eralisant l'utilisation faite dans (\EGA III~\oldS\S 4 et~5) du th\'eor\`eme de finitude pour un morphisme propre.

On applique ces derniers r\'esultats dans~\ref{X} et~\ref{XI}, consacr\'es respectivement \`a des th\'eor\`emes du type Lefschetz pour le groupe fondamental, et pour le groupe de \hbox{Picard}. Ces th\'eor\`emes consistent \`a comparer sous certaines conditions les invariants ($\pi_1$ ou $\Pic$) attach\'es respectivement \`a un sch\'ema $X$ et \`a un sous-sch\'ema $Y$ (jouant le r\^{o}le d'une section hyperplane), et \`a donner notamment des conditions o\`u ils sont isomorphes. Grosso modo, les hypoth\`eses faites servent \`a passer de $Y$ au compl\'et\'e formel de~$X$ le long de~$Y$, et \`a pouvoir appliquer ensuite les r\'esultats de~\ref{IX} pour passer de l\`a \`a un \emph{voisinage} ouvert $U$ de $Y$ dans $X$. Pour pouvoir passer de $U$ \`a~$X$, il faut disposer encore de renseignements (type \og puret\'e\fg ou \og parafactorialit\'e{\fg}) pour les anneaux locaux de $X$ en les points de $Z=X-U$, (qui est un ensemble fini discret dans les cas envisag\'es). Ceci explique l'interaction dans les d\'emonstrations des expos\'es~\ref{X}, \ref{XI}, \ref{XII} entre les r\'esultats locaux et globaux, notamment dans certaines r\'ecurrences. Les r\'esultats principaux obtenus dans~\ref{X} et~\ref{XI} sont les th\'eor\`emes \emph{de nature locale} \ref{X}~\ref{X.3.4} (\emph{théorème de pureté}) et \ref{XI}~\ref{XI.3.14} (\emph{théorème de parafactorialité}). On notera que ces th\'eor\`emes sont d\'emontr\'es par des techniques cohomologiques, de nature essentiellement globale. Dans~\ref{XII} on obtient en utilisant les r\'esultats locaux pr\'ec\'edents, les variantes globales de ces r\'esultats pour des sch\'emas projectifs sur un corps, ou plus g\'en\'eralement sur un sch\'ema de base plus ou moins quelconque; parmi les \'enonc\'es typiques, signalons \ref{XII}~\ref{XII.3.5} et \ref{XII}~\ref{XII.3.7}.

Dans~\ref{XIII}, nous passons en revue quelques uns des nombreux probl\`emes et conjectures sugg\'er\'es par les r\'esultats et m\'ethodes du S\'eminaire. Les plus int\'eressants peut-\^etre concernent les th\'eor\`emes du type Lefschetz cohomologiques et homotopiques pour les espaces analytiques complexes, cf. \ref{XIII} pages 26 et suivantes. Dans le contexte de la cohomologie \'etale des sch\'emas, les conjectures correspondantes sont prouv\'ees dans~\ref{XIV} par une technique de dualit\'e qui devrait s'appliquer \'egalement dans le cas analytique complexe (cf. commentaires \ref{XIII} p.\,25 et \ref{XIV}~\ref{XIV.6.4}). Mais les \'enonc\'es homotopiques correspondants dans le cas des espaces analytiques (et plus particuli\`erement les \'enonc\'es faisant intervenir le groupe fondamental) semblent exiger des techniques enti\`erement nouvelles (cf. \ref{XIV}~\ref{XIV.6.4}).

Je suis heureux de remercier tous ceux qui, \`a des titres divers, ont aid\'e \`a la parution du pr\'esent volume, dont les collaborateurs d\'ej\`a cit\'es dans cette Introduction. Plus particuli\`erement, je tiens \`a remercier Mlle CHARDON pour la bonne gr\^ace avec laquelle elle s'est acquitt\'ee de la t\^ache ingrate que constitue la pr\'eparation mat\'erielle du manuscrit d\'efinitif pour la photo-offset.

\bigskip\hfill Bures-sur-Yvette, Avril~1968
\par\smallskip \hfill A.\, Grothendieck.

